% ============================================================================
% Fichier  : partie0/P00_decouverte_rancongiciel.tex
% Titre    : Prologue : la Découverte
% Auteur   : MINKA MI NGUIDJOÏ Thierry Emmanuel
% ============================================================================

\chapter{Prologue : la Découverte}
\label{chap:prologue}
\chaptermark{Prologue}

\epigraph{%
    « Quand l'écran de rançon s'affiche, l'attaque est déjà vieille~;\\
    la seule question est de savoir si quelqu'un l'a vue venir. »%
}{Principe du cours}

\niveauChapitre{Introductif}{Aucun prérequis. Les notions sont reprises en
détail dans les chapitres.}

\begin{boxobjectifs}
\begin{itemize}
    \item Situer la supervision, le forensique et la réponse à incident dans
          une même journée réelle.
    \item Distinguer ce qui relève de la détection, de l'investigation et de
          la décision.
    \item Découvrir le fil rouge LiZbiZ et le Trilemme CRO.
    \item Formuler les questions auxquelles le cours répond.
\end{itemize}
\end{boxobjectifs}

\bigskip

% ============================================================================
\section*{6 mars 2026, 6 h 40, site de LiZbiZ à Garoua}
% ============================================================================

Le vigile de nuit a cru à une coupure~: les écrans de l'entrepôt affichent
tous le même fond noir et un texte en anglais réclamant un paiement. À
6~h~40, l'astreinte informatique est jointe. À 7~h~15, M.~KOUNOU, responsable
informatique, constate que le serveur de fichiers \texttt{FS-01}, la base de
suivi des colis et deux postes de la comptabilité sont chiffrés. Les
sauvegardes de la semaine sont, elles aussi, inaccessibles.

\medskip
Ce que personne ne sait encore ce matin-là\footnote{Cas fictif établi pour le
cours. Les noms (LiZbiZ, NYONG SÉCURITÉ, BANQUE SANAGA), les dates et les
montants sont fictifs~; les mécanismes techniques et les références
normatives sont exacts.}~: l'attaque n'a pas commencé cette nuit. Les
premières traces, dans les journaux, remontent à onze jours plus tôt. Entre
les deux, un compte de service compromis, un déplacement discret d'une
machine à l'autre, une exfiltration de 40~Go, puis le chiffrement. Pendant
onze jours, tout était dans les journaux. Personne ne les lisait.

% ============================================================================
\section{Trois métiers, une seule journée}
\label{sec:p00-metiers}
% ============================================================================

Cette journée met en scène les trois compétences du cours, dans l'ordre où
elles interviennent.

\begin{tableaucours}{L{3.2cm}L{4.4cm}L{4.6cm}}{%
    \tableauHeader{Moment} &
    \tableauHeader{Compétence} &
    \tableauHeader{Question posée}
}
Avant (onze jours) & Supervision, détection & L'a-t-on vu venir~? \\
Pendant (le constat) & Réponse à incident & Comment limiter les dégâts~? \\
Après (l'enquête) & Analyse forensique & Que s'est-il exactement passé, et peut-on le prouver~? \\
\end{tableaucours}
\vspace{-0.8em}
\begin{center}{\small\itshape Tableau~0.1. Les trois temps d'un incident}\end{center}

\begin{boxscenario}{Opération LOGONE. La journée du 6 mars}
\textbf{Constatation.} À 7~h~15, quatre systèmes chiffrés, sauvegardes
hebdomadaires inaccessibles, un message de rançon. La production de
l'entrepôt de Garoua est arrêtée.

\medskip
\textbf{Analyse} (confiance moyenne au matin, élevée après enquête).
L'incident n'est pas un simple chiffrement~: c'est l'aboutissement d'une
intrusion de onze jours. La distinction entre ce que l'on \emph{constate}
(les écrans noirs) et ce que l'on \emph{suppose} puis \emph{établit} (le
chemin de l'attaquant) est le cœur du métier.
\end{boxscenario}

% ============================================================================
\section{Le fil rouge : LiZbiZ et ses opérations}
\label{sec:p00-filrouge}
% ============================================================================

Tout le cours suit la même entreprise, \sigle{LiZbiZ}, transporteur de colis
présent à Yaoundé, Douala et Garoua, et le cabinet \sigle{NYONG SÉCURITÉ}
qui l'assiste. Trois opérations, situées dans le Nord, structurent la
pratique~:

\begin{tableaucours}{L{3.0cm}L{9.2cm}}{%
    \tableauHeader{Opération} &
    \tableauHeader{Ce qu'elle illustre}
}
BÉNOUÉ & La supervision et la détection~: ce que le SIEM a vu, et pas vu, pendant les onze jours \\
FARO & L'analyse forensique~: disque, mémoire, réseau, et la preuve opposable \\
MAYO & La réponse à incident~: confinement, éradication, rétablissement, retour d'expérience \\
\end{tableaucours}
\vspace{-0.8em}
\begin{center}{\small\itshape Tableau~0.2. Les trois opérations du fil rouge}\end{center}

\begin{tableaucours}{L{3.4cm}L{8.8cm}}{%
    \tableauHeader{Acteur} &
    \tableauHeader{Rôle}
}
M.~KOUNOU & Responsable informatique de LiZbiZ \\
Mme~NKENG & Analyste SOC, cabinet NYONG SÉCURITÉ \\
M.~BIKORO & Analyste forensique et chef de mission \\
Mme~MOUAFO & Responsable de la réponse à incident \\
M.~ETEME & Directeur général de LiZbiZ \\
\end{tableaucours}
\vspace{-0.8em}
\begin{center}{\small\itshape Tableau~0.3. Les acteurs de l'opération LOGONE}\end{center}

% ============================================================================
\section{Le Trilemme CRO}
\label{sec:p00-cro}
% ============================================================================

Ce cours partage avec l'investigation numérique un fil conducteur théorique~:
toute trace, tout acte d'enquête se juge sur trois dimensions à la fois.

\begin{boxCRO}
\textbf{Le Trilemme CRO.} \termecle{Confidentialité}~: la trace et l'enquête
ne doivent pas exposer plus de données que nécessaire. \termecle{Fiabilité}~:
la trace est-elle fidèle, complète, non altérée~? \termecle{Opposabilité}~:
pourra-t-on s'en prévaloir devant un tiers (direction, ANTIC, juge)~?

\smallskip
Une action d'enquête améliore souvent une dimension au détriment d'une autre.
On la note \cro{+}{-}{0}~: ici, confidentialité préservée, fiabilité
dégradée, opposabilité inchangée. Le professionnel arbitre en conscience et
\emph{documente} son arbitrage.
\end{boxCRO}

% ============================================================================
\section{Les questions du cours}
\label{sec:p00-questions}
% ============================================================================

\begin{boxinfo}
\textbf{Ce que le cours apprend à répondre.}
\begin{itemize}
    \item Comment \textbf{superviser} un système pour voir l'attaque pendant
          qu'elle se déroule (Parties~I et~II)~?
    \item Que disent les \textbf{normes et le droit} sur la preuve et le
          signalement (Partie~III)~?
    \item Comment \textbf{investiguer} proprement disque, mémoire et réseau
          (Partie~IV)~?
    \item Comment \textbf{répondre} à l'incident et s'en relever (Partie~V)~?
    \item Comment \textbf{apprendre} et mesurer, pour que la prochaine fois
          quelqu'un voie venir l'attaque (Partie~VI)~?
\end{itemize}
\end{boxinfo}

\separateur

\begin{boxresume}
\textbf{Ce qu'il faut retenir de ce prologue~:}
\begin{itemize}
    \item \textbf{L'écran de rançon} est la fin d'une histoire, pas le début.
    \item \textbf{Trois temps}~: voir venir (supervision), limiter (réponse),
          établir (forensique).
    \item \textbf{Le fil rouge} LiZbiZ et les opérations BÉNOUÉ, FARO, MAYO.
    \item \textbf{Le Trilemme CRO}~: confidentialité, fiabilité, opposabilité.
    \item \textbf{Constater} n'est pas \textbf{supposer}~; on sépare toujours
          les deux.
\end{itemize}
\end{boxresume}

\begin{boxinfo}
\textbf{Transition.} Pour voir venir l'attaque, il faut d'abord comprendre
pourquoi et comment on supervise~: c'est l'objet du
Chapitre~\ref{chap:pourquoi}.
\end{boxinfo}
